\chapter{Interval Estimations}
\label{chap:intervals}

\vspace*{2cm}

\begin{center}
	\itshape
	“Misinterpretation and abuse of statistical
	tests, confidence intervals, and statistical
	power have been decried for decades, yet
	remain rampant. A key problem is that there
	are no interpretations of these concepts
	that are at once simple, intuitive, correct"
	
	\vspace{0.5em}
	The ASA's Statement on p-Values: Context, Process, and Purpose
	Wasserstein, R. L. and Lazar, N. A., The American Statistician
	70, 129 (2016)\todo{Add ref}
\end{center}

\vspace{2cm}

In our eternal fight to try to infer something about the unknown parameters that govern the observables we measure, we introduce a new type of inference: the interval estimation. 

We are looking for a way to express, up to a certain degree of \emph{confidence}\sidenote{We'll define what confidence really is in a moment, so just bear with us} what interval the parameter that we are trying to estimate falls within. We take a second to note that while it may seem that we have already achieved something similar, this request is fundamentally different than anything our point estimate theory can provide. A naïve reader may say:

	\[
	\hat\theta = a \pm \sigma_a \Rightarrow \hat\theta \in \left[a-\sigma_a, a+\sigma_a\right]
	\]

Though you, as an attentive reader and excellent statistician will immediately identify the crucial flaw in this statement: $\sigma_a$ indicates the \emph{variability} of the estimate in the ensemble of repeated experiments, and does not represent the same "hard" upper/lower bound that we are trying to achieve through interval estimations. Really, the values $a-\sigma_a, a+\sigma_a$ have no more significance than any others in the point estimate formulation.

\begin{example}
	If we take a second, we can even find some edge case examples where the point estimate method breaks down, considering a Poisson pdf. We know that the MLE is $\hat\mu=k$ with null bias and $\Var\left[\hat\mu\right] = \mu$. Therefore if I measure one count $k=1$ then I obtain the point estimate for $\mu = 1\pm1$. Only armed with this we would be tempted to say that we can't exclude $\mu\leq0$, even though we know that the counts can, by definition, only be positive.
	
	Even more pathologically, if we measure no counts, we arrive to $\mu=0\pm0$, which, taken at face value, means that we have an infinitely precise estimate of the mean and that it is null. 
	
	Of course, this example shows a breaking point for our theory up until this point, and it invites us to explore new methods particularly useful for cases such as this one.
\end{example}
\todo{Low-n Poisson histogram graph? CBTM (could be the move)}

When we talked about point estimations, we never came to talk about intervals in which the parameter was confined, and going to an arbitrarily low probability basically any value was permitted. Now, we want to formulate a theory that limits, up to a certain degree of confidence, what values the parameter can adopt. 

\begin{definition}[Interval estimator]
	Given a pdf $p(x;\theta)$ dependent on an unknown parameter $\theta\in A$, an \emph{interval estimator} is a statistic $f(x): X\rightarrowtail \mathcal{P}(A)$ which outputs a subset of the parameter space.
\end{definition}

While a point estimate provides a single "best guess," an interval estimate offers a quantitative measure of uncertainty. At first glance, reporting a range might seem less precise than inferring a specific value, but in practice, they serve distinct and vital roles. Interval estimates are, but not limited to, merely a contingency for when point estimates become unstable, such as in low-count Poisson regimes, but are essential tools for experimental design and analysis. They allow us, for example, to determine the required duration for data collection to achieve a desired precision. Furthermore, they are often the only rigorous way to report results in "null" observations, such as setting an upper limit on the probability of a rare event that has yet to be measured.

As we have mentioned many times during this course, it is the job of the expert data analyst to know what tools to apply to a certain data set in the specific situation they might find themselves in. 

\section{Bayesian interval estimation}
We now consider the "fortunate" case in which the unknown parameter can be described by a probability density function, allowing us to perform inference using Bayes' theorem, as discussed in Section \ref{subsec:Inf_Bayes}. Starting from Eq. (\ref{eq:posterior_dob}), we write the posterior distribution as
	
	\[
	\pi_x(\theta) = \pi(\theta)\frac{L_x(\theta)}{p(x)}
	\]

We then determine a subset $f_{x}(\theta)$ of the parameter space\sidenote{For this chapter, we will adopt this notation, where the quantity that is fixed will be the subscript and the quantity that represents the domain will be in parentheses. In this case, $f_{x}(\theta)$ is a subset of the parameter space, at a fixed $x$} (typically an interval), depending on the observed data $x$, such that
\begin{equation}
	\int_{f_{x}(\theta)} \pi_x(\theta)\, d\theta = c
	\label{eq:Bayes_int_eq}
\end{equation}
where $0 < c < 1$ is a constant called the \emph{credibility}. The set $f_{x}(\theta)$ represents a \emph{credible region}, i.e. a set of plausible values for $\theta$ such that the posterior probability of $\theta$ lying within it is equal to $c$:
	
	\[
	P(\theta \in f_{x}(\theta)\mid x) = c
	\]

The importance of this scripture cannot be overstated: we are finding an interval in the \emph{parameter space}, and this interval \emph{depends on the observed data}.

The credibility $c$ quantifies the degree of belief that the true value of the parameter lies within the region $f_{x}(\theta)$. In practice, $c$ is chosen in advance, typically taking values between $0.9$ and $0.95$.

As usual, the Bayesian formulation brings many nice properties such as respecting the Likelihood principle, obtaining clear, specific, and intuitive results. However, it depends on how we choose the prior, which can greatly influence the type of result we obtain. Also, we need $\theta$ to be able to admit a pdf in the first place. \todo{Here we can talk about Jeffrey's prior, if it ever comes up, or reference priors.}

Integral (\ref{eq:Bayes_int_eq}) admits infinite solutions, so we must introduce a criteria to be able to select the specific solution we want. The most agreed upon way to proceed is by selecting an \emph{ordering function}\sidenote{We will dedicate more time to ordering functions in the near future} $o_{x}(\theta)$ that assigns a priority hierarchy to the values of the parameter to include. 

In this case, the equation becomes:
	
	\[
	\int_{o_{x}(\theta)\geq q_c} \pi_x(\theta)\dd \theta = c
	\]

Where $q(c)$ is a constant that depends on the credibility $c$ in a way such as to satisfy the equation. Usually, this allows us to find one unique solution, given $c$ and the ordering function. Essentially, we are finding an interval in $\theta$ where $o_{x}(\theta)\geq q_c$ and the above integral is equal to $c$.

The most common and simple choice is to use \emph{posterior ordering} $o_{x}(\theta) =\pi_x(\theta)$, where the $\theta$ that greater contribute to $\pi_{x}(\theta)$ are integrated over first. Essentially, we want to prioritize "credible" $\theta$ values and so we choose the integral limits so that those are included first. This may seem complicated, but it's nothing more than integrating the posterior starting from its maximum and expanding "outwards", until the area underneath the curve equals the credibility. This algorithm certainly will output a credible region containing the most credible $\theta$, as well as the one with the smallest measure in $\dd\theta$. This is desirable, as it intuitively restricts the credible values of the parameters more than any other ordering function. We note that this is not invariant for variable changes in the parameter space. 

We will talk about ordering algorithms in greater detail in section \ref{subsec:ordering}.

\begin{example}
	Let's formulate an example: let's imagine to measure an observable that is distributed following a Poisson pdf depending on a parameter $\mu$: $p(k;\mu) = \frac{e^{-\mu}\mu^k}{k!}$. However, let's assume that for some reason we know that $\mu$ admits a pdf and that it is uniformly distributed between $0$ and $10$. 
	
	I can therefore write my posterior:
	
	\[
	p_k(\mu) = \underbrace{\frac{1}{10}}_{\text{prior}}\overbrace{\frac{e^{-\mu}\mu^k}{k!}}^{\text{Likelihood}}\cdot\frac{1}{p(k)}\quad 0\leq\mu\leq10
	\]
	
	As far as $p(k)$ is concerned, it is the value that normalizes the posterior. In other words, it is the marginal probability of observing k under our whole prior model:
	
	\[
	p(k) = \int_0^{10} \frac{e^{-\mu}\mu^k}{k!}\cdot\left(\frac{1}{10}\right)\dd\mu
	\]
	
	At the end of the day, this complex expression is no more than a multiplicative factor that makes it so that the posterior $p_k(\mu)$ is correctly normalized, so we will just leave it written as $p(k)$ knowing this.
	
	The following figure shows the posterior for different values of $k$, and as we expect we see that for each $k$ the function peaks around the observed value.
	
	\begin{figure}[H]
		\centering
		\input{images/bayes_posterior.pgf}
		\caption{Posterior for different values of $k$ in the $\mu$ space}
		\label{fig:bayes_posterior}
	\end{figure}
	
	We want to find, given a measurement $k$, an interval for $\mu$ with a credibility of $90\%$. We know this to mean that we want to find an interval in $\mu$ of the type $\mu\in[a,b]$ such that $\int_a^bp_{k}(\mu)\dd \mu= 0.9$. We can imagine there to be an infinite number of such intervals, and so we define different \emph{ordering functions} that change the type of interval we get.
	
	For example, if we choose to integrate over the zone where $p_{k}(\mu)\geq q_c$, we define two bounds (the two solutions $a,b$ of $p_{k}(\mu)\geq q_c$) such that 
	
	\[
	\int_{p_{k}(\mu)\geq q_c}p_{k}(\mu)\dd \mu = 0.9
	\]
	
	This is called "posterior ordering" and is the most natural choice, as it includes the maximum value of the posterior by definition, meaning that it also results in the minimal interval length in $\mu$.
	
	We can also integrate by starting from the maximum or minimum value of $\mu$, obtaining lower or upper estimates of the parameter, respectively. 
	
	These different orderings are shown in the following image, assuming to have observed $k=5$.
	
	\begin{figure}[H]
		\centering
		\input{images/bayes_interval.pgf}
		\caption{Different ordering intervals for $k=5$.}
		\label{fig:bayes_interval}
	\end{figure}
	
	\input{tables/bayes_interval.tex}
	
	As we can see, the posterior ordering results in the smallest interval. One may be tempted to combine the lower and upper bounds into a smaller interval than even the posterior ordering, but doing so would break the definition of credibility (not only in the mathematical term, but also your credibility as a data analyst), as we would be artificially shortening an interval and therefore lowering the credibility.
	
	One may also be tempted to swap between lower and upper limits depending on the measured value of $k$. However, this temptation must be resisted, and in the future we will go into detail as to why this habit is to be avoided (\emph{flip-flopping} between ordering functions).
\end{example}

\todo{Add to the example the variable change stuff?} 

We can see the power of this Bayesian case, and it's not hard to understand what an interval estimate looks like and means under this formulation. However, it is not always applicable, as the unknown parameters must admit \emph{pdf's} themselves, and we must know or assign them, risking the loss of generality.

\section{The standard interval estimator}

We once again move away from the Bayesian outlook, which remains intuitive and extremely descriptive, though limited by the assumption of a \emph{prior} $\pi(\theta)$.

In the general case, in order to find an interval estimate for a certain parameter $\theta$ we will have to solve the following equation, which synthetically gives us the \emph{operative definition of a confidence interval}:

\begin{equation}
	\int_{o_{\theta}(x)\geq q(CL)} p(x;\theta)\dd x \geq CL
	\label{eq:int_definition}
\end{equation}

Let's look at it generally with the idea to grasp the basic concept, then in the following sections we will elaborate on each part individually.

Basically, at its core, we can condense this entire equation into the most basic form of:

	\[
	g_{x_0}(\theta) - CL \geq 0
	\]

If we think about it, on the left side we have an integral in $\dd x$, and on the right a constant that we fix beforehand, namely the "confidence level" \emph{CL}\sidenote{Don't worry about these definitions, they will all get cleared up in the near future. This part aims to provide a general outlook on the math that we will be using, so just bear with us}. Fundamentally, solving the inequality at a fixed $x_0$ dependent on the parameter returns an interval (or family of disconnected intervals) for which the inequality is true. These $\theta$ intervals are the ones we will be studying.

Now, what exactly is this $g_{x_0}(\theta)$? Well, since we are not in the Bayesian case and cannot directly say that the parameter admits a pdf, we can only talk about the distribution of the observable: we want that, with probability \emph{CL}, the interval defined by $g_{x_0}(\theta) - CL \geq 0$ contains $\theta$, to be intended in the frequentist sense. 

As we will see, infinite such intervals exist for each fixed $\theta$ (there are infinite ways to integrate a normalized function such that its integral is equal to, say, $0.95$), and $o_{\theta}(x)$ serves as an "ordering function" with the aim of unambiguously defining one such interval by describing which $x$'s to include. 

This integral in eq. (\ref{eq:int_definition}) is nothing more than the integral of the pdf itself, which we know to be the cumulative of the density function\sidenote{This integral can be seen as, at most, the difference of two \emph{CDF's} in order to make the bounds work out}, calculated on upper and lower bounds that may, in general, depend on the parameter. 

What this all means is that having fixed an observed $x_0$ (coming from pdf $p(x;\theta)$), \textbf{the ordering function describes a $\theta$-dependent $x$-interval over which the integral of the pdf is greater than or equal to the \emph{CL}}. Therefore, after having measured $x_0$, we can solve this expression (which we so naïvely wrote as $g_{x_0}(\theta) - CL \geq 0$) to find an interval in $\theta$, and that will be the estimate that we are looking for. 

Solving this equation is rarely simple, and can be done when the pdf assumes a closed form cumulative, or when the pdf is written using pivotal quantities so that its cumulative is tabled. The next sections will break down each part of the process, using examples and visual aids to try to better explain this difficult, yet fundamental concept.

\subsection{Confidence intervals}
Let's begin our piece-by-piece breakdown of eq. (\ref{eq:int_definition}) by really understanding what our goal here is. For a given observation $x_0$ coming from $p(x;\theta)$, we wish to write an interval in $\theta$ that contains the true value $\theta_t$ of the parameter with a probability of \emph{CL}. In other words, a specific confidence interval is a member of a set, such that the set has the property that\cite{Feldman_1998}:

\begin{equation}
	P\left(\theta\in\left[\theta_1,\theta_2\right]|x_0\right) = CL
\label{eq:confidence_interval}
\end{equation}

Here, $\theta_1$ and $\theta_2$ are functions of $x_0$, and are therefore different for each measurement of the observable, even at a fixed $\theta$. For a fixed set of confidence intervals (i.e, for a fixed $x_0$), we wish for the equation to be true for all allowed $\theta$. This means that we can interpret the above statement to be that the interval contains the \emph{fixed, unknown, true} $\theta_t$ in a fraction \emph{CL} of experiments.\sidenote{This is different than the Bayesian case where we could state that the \emph{degree of belief} that $\theta_t\in\left[\theta_1,\theta_2\right]$ is \emph{CL}}

If equation (\ref{eq:confidence_interval}) is satisfied for all $\theta$, then one says that the intervals "cover" $\theta$ at the stated confidence. For any values of $\theta$ in which the equation is not satisfied we say that the interval "undercovers" or "overcovers" for that $\theta$, depending if the probability is less than or greater than \emph{CL}. We naturally don't want undercoverage, and reluctantly accept overcoverage as an unfortunate price to pay in the precision with which we can construct the interval. Generally, a set of intervals for which there is no undercoverage for any $\theta$ and some overcoverage for some $\theta$ is considered "conservative", and while not the end of the world, means that these intervals will generally be wider and less informative on the true value of $\theta$ than the confidence level quoted.

\begin{figure}[H]
	\centering
	\input{images/binomial_coverage.pgf}
	\caption{3 different coverage algorithms for $20$ samples of a binomial distribution}
	\label{fig:binomial_coverage}
\end{figure}

Figure \ref{fig:binomial_coverage} demonstrates the conservatism inherent in discrete interval estimation. Because the number of successes in 20 trials must be an integer, the coverage probability $\mathcal{C}(p)$ cannot be a flat line at $0.95$; instead, it forms a "sawtooth" pattern. Note that all three algorithms satisfy the definition of a $95\%$ confidence level because their global minimum ($\inf_p \mathcal{C}(p)$) is $\ge 0.95$. The probability ordering is generally preferred\sidenote{Well, for now at least, since it is the direct parent of the posterior ordering seen in the Bayesian case. Later, we'll explore the more general two-sided \emph{Likelihood ratio} ordering} for two-sided intervals as it minimizes over-coverage, staying closer to the nominal level. In contrast, the Lower and Upper orderings are essential when one is specifically interested in one-sided bounds, despite their extreme conservatism near $p=0$ and $p=1$.

From this whole part we want to take away one very important concept: when we obtain an interval $[a,b]$ with confidence level \emph{CL}, that does not mean that the true value of the parameter $\theta_t$ lies within that interval with probability \emph{CL}. What it means, instead, is that the interval was generated by a procedure whose repeated use contains the true parameter with probability \emph{CL}. Before moving on, this concept should be clear, as it highlights the key difference between the Bayesian and frequentist approaches to interval estimations.

\subsection{Ordering algorithms}
\label{subsec:ordering}

We have already briefly touched upon ordering algorithms up until this point, and now we'll expand on what we have already introduced and complete our arsenal of ordering functions with a few more neat algorithms.

The idea is to study the different options we have for $o_{\theta}(x)$ in the usual equation (\ref{eq:int_definition}), with the goal of removing the ambiguity in the way we choose the interval that solves it. What the ordering function is really doing is, for each fixed $\theta$, "ranking" (well, \emph{ordering}) each possible $x$ according to the value of the ordering function $o_{\theta}(x)$. From this, we add $x$ values to the confidence region for that $\theta$ based on this ordering until the sum of probabilities reaches the \emph{CL} to create an interval in $x$ for a fixed $\theta$.

Later, this will be repeated \emph{using the same ordering algorithm} for all $\theta$, allowing us to obtain the confidence interval for an observed $x_0$ as the set of all $\theta$ values whose acceptance regions contain $x_0$:

	\[
	I_{x_0}\left(\theta\right) = \theta: \{x_1(o_{\theta})\leq x_0 \leq x_2(o_{\theta})\} 
	\]

This shows that the $x_0$-dependent $\theta$ interval (the interval estimate that we are looking for) is defined as the set of all $\theta$ such that the observed $x_0\in\left[x_1,x_2\right]$, where $x_1,x_2$ are calculated through the ordering function at a fixed $\theta$. This expression is nothing different than what we have already stated before, and can be considered to be a restatement of eq. (\ref{eq:int_definition}), with the goal of repeating the same concept in different ways to aid in understanding what interval estimations are. CITE

Basically, as we have said before, solving the ordering function equation $o_{\theta}(x)\geq q(CL)$ we are finding the $\theta$-dependent integral bounds that meet the desired confidence level, with the desired ordering. Let's look at some common choices for ordering functions.

We could start with setting $o_{\theta}(x)=p(x;\theta)$ (\textbf{probability ordering}), which is a universal choice that brings the smallest interval in $x$ (which does not necessarily translate to the smallest interval in $\theta$). This means that we are ordering using the pdf itself, choosing the most probable values of $x$ in decreasing order. This should be reminiscent of what we did in the Bayesian case, where we were fortunate enough to be able to order through the posterior itself.

Another natural choice would be to use $o_{\theta}(x)=\pm x$, as to start from the lowest value of $x$ and include higher values until the confidence level is reached (essentially finding the value for which the cumulative is equal to \emph{CL}). We call this \textbf{lower ordering}, which brings us \emph{lower limits on} $\theta$ (one-sided intervals of the type $\theta\geq\theta_1$). 

	\[
	F_{\theta}(x_{max}) = CL \quad \Rightarrow \quad \theta\geq\theta_1
	\]

We could also do the inverse, as in starting from the highest value of $x$ and including lower values until the integral equals \emph{CL} (essentially finding the value for which the cumulative is equal to $1-CL$). We call this \textbf{upper ordering}, which brings us \emph{upper limits on} $\theta$ (one-sided intervals of the type $\theta\leq\theta_2$).

	\[
	F_{\theta}(x_{max}) = 1 - CL \quad \Rightarrow \quad \theta\leq\theta_2
	\]

We could also think to take the central band (\textbf{central ordering}), as to exclude $\frac{1-CL}{2}$ from the top and bottom. The ordering function $o_{\theta}(x)$ is defined implicitly by the tails of the distribution. We include $x$ values in our acceptance set starting from the median and expanding outwards until the probability excluded in each tail separately does not exceed $\frac{1-CL}{2}$, or alternatively repeating the upper/lower ordering simultaneously on the tails until half of \emph{CL} is reached. 

The final ordering algorithm we will explore for now is one based on the \textbf{Likelihood Ratio} (some texts may call it "Feldman-Cousins ordering"). The idea is to use the Likelihood Ratio as the ordering function
	
	\[
	LR_{\theta}(x) = \frac{p\left(x \middle| \theta\right)}{p\left(x \middle| \hat{\theta}_{MLE}\right)} = \frac{L_{x}\left(\theta\right)}{L_{x}\left(\hat{\theta}_{MLE}\right)}
	\]

We have written both versions, explicitly in $p$ and in $L$, and while they represent the same exact quantity both can be found in literature. The important thing to note is that the Likelihood ratio is a \emph{statistic}, that is, it is $x$-dependent quoted at a fixed $\theta$.

Though we will rarely follow this line of reasoning\sidenote{Usually we can only solve this numerically, and if we have high statistics we invoke Wilks' theorem, which we will talk about soon}, we can think of this ordering algorithm as the solution to the following system of equations:
	
	\[
	\begin{cases}
		\int_{x_1}^{x_2} p(x;\theta)\dd x = CL \\
		LR(x_1) = LR(x_2)
	\end{cases}
	\]

The Likelihood ratio, which we will see in numerous forms\sidenote{Confusingly enough, they all go under the same name and while there are strong similarities between all instances of the LR, it's important to understand what our goal is and to not blindly start dividing Likelihoods.} serves as a statistic that compares some general-valued Likelihood with a specified (fixed parameter) Likelihood. In this case, for each fixed $\theta$ we are comparing the likelihood of $\theta$ with respect to the maximally likely one.

In terms of ordering, what this corresponds to is starting at the top of the $LR(x)$, and increasing the $x$ interval according to the $LR$ until the interval defined through $LR$ corresponds to an integral of measure \emph{CL} in $p(x)$. Conceptually, it's the same as probability ordering, where instead of ordering using $p(x)$ you order using $LR(x)$. This is where the above system comes from, we set a threshold on the Likelihood ratio such that the $x$ values of the intersection define the bounds of the confidence interval in $p$. This is favorable with respect to probability ordering, as the interval takes into consideration the information we have on $\theta$ (like its physical limits), and is guaranteed to contain the \emph{MLE}, as well as bring the smallest possible interval in $\theta$. Furthermore, by definition the interval will be fully covered for all $\theta$ in the continuous case, and will be, at worst, conservative in the discrete case.

\begin{figure}[H]
	\centering
	\input{images/lr_ordering.pgf}
	\caption{Visual representation of the Likelihood Ratio ordering algorithm. $x$ is distributed following a Normal distribution with unitary variance and unknown $\mu\geq 0$, where the case $\mu=0.5$ has been shown. The \emph{LR} and threshold have been translated upwards by $0.5$ for clarity, and the $95\%$ confidence region for $p(x)$ is shown.}
	\label{fig:lr_ordering}
\end{figure}

Figure \ref{fig:lr_ordering} gives a visual representation of this method and how to use it to find the bounds in $x$ for a specific value of $\mu$. Within it, the condition that $\mu$ must be greater than zero is considered, which brings the asymmetric tail towards negative $x$ which probability ordering would not have accounted for. Naturally, this image is for a general fixed $\mu$, and the left and right bounds depend on this value.

\subsubsection{Neyman Construction}
A way to visualize the effect of these different ordering algorithms is through the construction of a "belt" that associates, given the ordering algorithm, each measured value of the observable to an interval estimate. This plot is namely the \emph{Neyman Construction}, is built by solving eq (\ref{eq:int_definition}) for all fixed values of $\theta$. 

Given an observation of $x$, the interval in $\theta$ is found by intersecting the horizontal line $x=x_{obs}$ with the constructed belt, and taking the two points of intersection as the bounds for the parameter. The following image shows the three Neyman belts constructed from the previous example of $20$ extractions of a binomial with unknown success parameter $p$, using three different orderings. 

\begin{figure}[H]
	\centering
	\input{images/binomial_Neyman_belts.pgf}
	\caption{Neyman belts built for three different ordering algorithms for $20$ samples of a binomial distribution}
	\label{fig:binomial_Neyman_belts}
\end{figure}

For any outcome $x\in\{0,1,2,\dots,20\}$ we can draw a horizontal line in any of the belts in figure \ref{fig:binomial_Neyman_belts} and find our interval estimation for $p$. As we have already touched upon briefly, it's important to stick to one belt (and therefore ordering algorithm) for all measured values rather than \emph{flip-flopping} between them. If we "artificially" make the interval smaller by choosing a particular intersection or combination of different ordering algorithms, it only means that we have lowered the coverage below our \emph{CL}, and so we lose the very property that defined our intervals in the first place. The choice of ordering should be determined by the goal of the analyst, like choosing Likelihood Ratio for the inclusion of the \emph{MLE} or lower ordering for the upper bound.


\begin{example}
\label{ex:LR_vs_p_ordering}
	Let's build the Neyman band for a Normally distributed $x$ with unknown parameter $\mu$ (let's say $\sigma=1$ for simplicity). We choose, for the sake of this example, to use probability-ordering. We can't integrate the Normal distribution analytically, so we'll have to fall back on the tables found in Appendix \ref{tables:normal}.
	
	Specifically, what we want 
	
	\begin{figure}[H]
		\centering
		\input{images/neyman_gauss.pgf}
		\caption{Neyman belt for a Normally distributed $x$ with fixed unitary variance. Here, we have used probability ordering and we note that an interval exists for all possible values of $x$.}
		\label{fig:neyman_gauss}
	\end{figure}
	
	As seen in \ref{fig:neyman_gauss}, the belt is of simple interpretation: having fixed the variance the interval estimate is centered on the measured value with a defined width, and such an interval exists for all measured $x$. However, let us restrict ourselves in the particular case of $\mu>0$. In this case, if we simply use probability ordering again, it would be analogous to truncating the above belt, excluding all negative $\mu$ values.
	
	However, doing this we realize that not all observable values of $x$ correspond to a possible interval in $\mu$, as the following plot shows:
	
	\begin{figure}[H]
		\centering
		\input{images/neyman_gauss_limited.pgf}
		\caption{Neyman belts for a Normally distributed $x$ with fixed unitary variance and $\mu>0$. On the left, we have used probability ordering, and on the right Likelihood Ratio ordering. Plotted also is the line for an observation at $x=-2.5$}
		\label{fig:neyman_gauss_limited}
	\end{figure}
	
	Figure \ref{fig:neyman_gauss_limited} shows the breaking point of the probability ordering, where it fails to find an interval for $\mu$ for a negative $x$. On the right, the same belt is constructed with the sole difference of using a different ordering function: the Likelihood Ratio. As we can see, this new ordering function once again allows us to correspond every measure of $x$ to an interval, and takes into consideration the limits on the parameter itself. 
	
	We can also see that the interval in $x$ defined for $\mu=0.5$ corresponds to what we have shown in figure \ref{fig:lr_ordering}.
\end{example}

These belts are of great didactic caliber, and are useful in visualizing what it is that we are actually doing under this whole mathematical abstraction. Numerically, we are representing the dependence of the $\dd x$ integral bounds in $\theta$, such that $\int_{}$

Analytically, the goal would be to define a function, or set of functions, of the type $x_{max}(\theta)$, $x_{min}(\theta)$, or their inverses. Sometimes, based on the ordering algorithm, pdf, and parameter we are trying to estimate, we are able to build the whole belt analytically and define an interval for $\theta$ for all observed values of $x$. 


\begin{example}
	Let's build a Neyman belt analytically! We will study the common case of an observable $x$ distributed following a uniform pdf $p(x;m) = \frac{1}{m}$ for $0\leq x \leq m$ and $0$ outside of that support. We observe one $x$ and wish to build an interval estimate for $m$ with a $90\%$ \emph{CL}. 
	
	Let's start by finding a lower limit for $m$, which means using "upper ordering":
	
	\[
	\int_0^{x_{max}} \frac{1}{m}\dd x = CL
	\]
	
	Remember, we integrate with a fixed $m$, and have set $x_{min}=0$ thanks to our ordering function. We can easily solve to find:
	
	\[
	x_{min} = 0 \quad x_{max} = m\cdot CL
	\]
	
	These two equations define a Neyman belt, from which we easily see that the one-sided interval for $m$ is $m\geq \frac{x_0}{CL}$, where $x_0$ is the result of our measured observable. This belt is seen on the left side of figure \ref{fig:uniform_belts}.
	
	Let's solve the same problem, this time using \emph{LR} ordering to obtain a two-sided interval for our parameter. We know that for a uniform distribution, as we have worked out in appendix \textcolor{red}{REF} the \emph{MLE} for $m$ is $m=x$ itself, and the \emph{LR} function can be written as:
	
	\[
	LR(x) = \frac{x}{m}
	\]
	
	To satisfy Likelihood Ratio ordering, we pick the x with the largest values of \emph{LR} first, and so our interval in $x$ will be $[x_{min}, m]$, chosen so that:
	
	\[
	\int_{x_{min}}^m \frac{1}{m}\dd x = CL
	\]
	
	Solving the above integral, we get that
	
	\[
	x_{min} = (1-CL)\cdot m \quad x_{max} = m
	\]
	
	which, after a measurement $x_0$, translates to an interval of $x_0\leq m \leq \frac{x_0}{1-CL}$. We also note that in this particular case, \emph{LR} ordering corresponds exactly to having used lower ordering (which gives us the upper bound, though the lower bound is fixed by the fact that $x\leq m$).
	
	The following figure shows the two different belts that we were able to construct analytically for this problem, and it's clear how inverting the $x$ limits as a function of $m$ brings us the interval for $m$.
	
		\begin{figure}[H]
		\centering
		\input{images/uniform_belts.pgf}
		\caption{Different Neyman belts for the interval estimate on $m$, the parameter for a uniform distribution with minimum point $0$ and maximum $m$. On the left, upper ordering is shown and on the right Likelihood Ratio ordering. An example measurement of $x_0=1$ is shown, to highlight how the belt can be inverted to obtain a one or two sided interval estimate.}
		\label{fig:uniform_belts}
	\end{figure}
	
	In this specific case with $90\%$ \emph{CL}, we use as an example the measurement $x_0=1$, which brings the intervals of:
		
		\input{tables/uniform_belts.tex}
		
	In this case, we are more interested in an upper limit for $m$ so Likelihood/lower ordering would bring more information on the parameter, as the observation of $x_0$ itself already provides a natural lower limit for the value of $m$. 
\end{example}

However, the belt is not always possible to express analytically, and so often we must use numerical techniques or searched tabular data to find the interval we desire. Naturally, not every possible parameter-pdf combination will be expressed in tables, and for that reason we must find a way to transform our specific cases into general, standard problems that have already been solved before us. Luckily, in the next section we will introduce \emph{pivotal quantities} and how they can be used to standardize certain types of problems.
\\
\todo{I've found that looking at the scripts generating the graphs is really useful to understand how the different ordering algorithms and belts works. It could be a good idea to show small parts of the code as an additional way of displaying such concepts (for those who knows their python, that I suspect are vast majority of the possible readers). If I'm not wrong Prof. Baldini did something similar in his iconic Lab 1 notes.}
\todo{A clean way to present this material might be to include references, where appropriate, to a dedicated appendix collecting the relevant code. For example: ``To gain further insight into this concept, see Appendix~X, where the corresponding implementation is provided.''}

\section{Pivotal quantities}
As we have said, determining confidence regions starting from their direct definition can be especially difficult, even numerically. One way to simplify the problem is to search for \textbf{pivotal quantities}, which we take a second to define:

\begin{definition}[Pivotal quantity]
	Given a pdf $p(x;\theta)$, a function $s(x,\theta)$ is said to be a \emph{pivot} if and only if $\forall \theta_1,\theta_2\in A$ we have that:
	
	\[
	p\left(s(x,\theta_1);\theta_1\right) = p\left(s(x,\theta_2);\theta_2\right) \equiv p_s(s)
	\]
\end{definition}

From this definition some confusion may remain as to what exactly we are looking for: we are trying to find a function\sidenote{Question: is such a function a statistic? Hopefully you answered no, as it is, in general, a function that depends on unknown parameters!} of our observable(s) and parameter(s) such that the pdf of that function does not depend on the unknown parameters. In a certain sense, we are finding a transformation such that we "standardize" our pdf to a more general form, independent on what we don't know. 

For example, consider a Normally distributed observable, coming from $\N(x;\mu,\sigma)$. Let's write:

	\[
	z\left(x,\mu,\sigma\right) = \frac{x-\mu}{\sigma}
	\]\sidenote{For this specific pivot the variable $z$ is most commonly used.}

Well, now to write out $p_z(z)$, we need to use the change of variable rule that we worked out way in definition \ref{def:pdf_var_change}, essentially substituting and multiplying by the Jacobian of the inverse transformation:

	\[
	p_z(z) = \N(x(z))\left|\frac{\dd x}{\dd z}\right| = \frac{1}{\sqrt{2\pi}}e^{-\frac{z^2}{2}}
	\]

We see that the above expression is nothing more than a "standard Gaussian", which is a Normal distribution with null mean and unitary variance. This is independent of our original parameters, and all of the information we could ever hope to have on this distribution can be found on tables\todo{Perhaps ref ours when we have them} and in literature, as it is an extensively researched function. 

A possible strategy for when we have to deal with Normal distributions could then be to transform them through the pivotal quantity $z$, solve whatever problem it is that we are looking to solve (for example, but not limited to, finding an interval estimate for a parameter), and then transform $z$ back to obtain the necessary information on our parameters. 

Let's look into a key property of pivotal quantities that is in line with what we have been talking about for the last few pages. Let's consider an interval estimate expressed like $f(x): \{\theta: g\left(s(x,\theta)\right)>0\}$\sidenote{My confidence interval is the set of all possible $\theta$ values that make my pivotal test statistic $s$ fall within a certain range defined by the ordering function $g$}. In this case, we have \emph{coverage}:
	
	\begin{align*}
		\mathcal{C} &= \int_{x\in f(x)} p(x;\theta)\dd x = \int_{g\left(s(x,\theta)\right)>0} p(x;\theta)\dd x \\
		&= \int_{g(s)>0} p_s\left(s(\theta);\theta\right)\dd s = \int_{g(s)> 0} p(s) \dd s
	\end{align*}

What does this mean? If we express our density function through the pivotal quantity then the coverage is constant in the parameter. That means that the confidence region is easily found for all $\theta$ as the solution to $g\left(s(x,\theta)\right)>0$. The pesky problem that we had, where we had to solve the integral for every possible fixed value of the parameter is now gone: once I solve it for one $\theta$, I have the full belt with constant coverage, pretty cool! 

We note that we have already explored an example of this kind, fig \ref{fig:neyman_gauss} shows the Neyman belt for a Normally distributed $x$ with unitary variance, which is exactly the belt that we would obtain for the pivot $z$\sidenote{We don't need for the observable to have unitary variance, we just use the pivot $z$ which corrects for whatever the variance is}. As we can see in that example, the belt is independent for $\mu$ and its width was found by consulting the tables for the central integral of a standard Gaussian.\todo{ref?}

Other pivotal quantities exist, and when possible it is always a good idea to abuse their general properties.


\subsection{Likelihood ratio theorem}

Not always are we able to easily identify pivotal quantities for our problems, and therefore it can be useful to build an approximate quantity that works in most cases.

\begin{theorem}[Likelihood Ratio (Wilks') theorem]
	Let $p(x;\theta)$ be a pdf with domain independent from $\theta\in \R^N$, and at least twice differentiable. In this case, the quantity:
	
	\begin{equation}
	\lambda\left(x,\theta\right) = 2\ln\left[\frac{ p(x;\hat{\theta}_{MLE})}{p(x;\theta)}\right]
	\tag{Log-Likelihood ratio}
	\end{equation}
	
	asymptotically has a universal distribution. This means that it is independent of the shape of $p(x;\theta)$ and real value of $\theta$ (only its dimension, otherwise known as \emph{degrees of freedom}.)
\end{theorem}

We note that using $\lambda(x,\theta)> q_{CL}$ as an ordering function is equivalent to using Likelihood Ratio ordering.

This easy to build quantity is a universal pivot for whichever pdf we start with, and we know its distribution $p(\lambda)$ is a $\chi^2_\nu$, which means that:

	\begin{equation}
	p_\nu(\lambda) = \frac{1}{2^{\frac{\nu}{2}} \Gamma(\frac{\nu}{2})}\lambda^{\frac{\nu}{2}-1}e^{-\frac{\lambda}{2}}
	\label{eq:chi_squared}
	\end{equation}

Asymptotically in the degrees of freedom, this distribution approaches a Gaussian with mean $\nu$ and variance $2\nu$. 

The proof of Wilks' theorem will not be discussed in this write-up, though those who are interested can take a look at the following example to see how we conclude that $\lambda$ is distributed like a $\chi^2$, and how it approaches a Gaussian for high degrees of freedom:

\subsubsection{Derivation of the asymptotic distribution}
Let's derive the asymptotic distribution of the $LLR$, using the fact that Wilks' theorem guarantees its generality. Therefore, we are free to start with whatever distribution we wish, and we decide to use $N$ samples of a Normally distributed variable. We write the log Likelihood:
	
	\begin{align*}
		\lambda(x,\mu) &= 2\ln\left(Ce^{-\sum_i\frac{(x_i-\bar{x})^2}{2\sigma^2}}\right) - 2\ln\left(C e^{-\sum_i \frac{(x_i-\mu)^2}{2\sigma^2}}\right)\\
		&= \frac{1}{\sigma^2}\sum_i\left[(x_i-\mu)^2-(x_i-\bar{x})^2\right] \\
		&= \frac{1}{\sigma^2}\sum_i\left[-2\mu x_i + \mu^2 +2\bar{x}x_i -\bar{x}^2\right] \\
		&= \frac{1}{\sigma^2}\left[-2\mu N\bar{x} + N\mu^2 +2N\bar{x}^2 -N\bar{x}^2\right] \\
		&= \frac{1}{\sigma^2}\left[-2\mu N\bar{x} + N\mu^2 +N\bar{x}^2\right]\\
		&= \frac{N\left(\bar{x} -\mu\right)^2}{\sigma^2} = \frac{\left(\bar{x}-\mu\right)^2}{\left(\frac{\sigma}{\sqrt{N}}\right)^2} \\
		&= \frac{\left(\bar{x} - \mu \right)^2}{\sigma^2(\bar{x})} = \left(\frac{\bar{x}-\mu}{\sigma(\bar{x})}\right)^2
	\end{align*}
	
As we can see, $\lambda$ is a pivot for all $N$, distributed like the square of a standard Gaussian variable. We can see that for an $n-$dimensional Gaussian with mean $\vec{\mu}$, we would have the sum of $n$ equal terms like the one above, one for each of the components of $\mu_i$. These universal distributions are called the \textbf{chi squared with $n$ degrees of freedom $\left(\chi^2_n\right)$}.
	
Let's start by writing $\chi^2_1$, which we get through squaring a Normally distributed $x$ and finding the pdf of this new variable:
	
	\begin{align*}
		p_y(y) &= \sum_i^{N_{sol}} p(x) \left|\frac{\dd x}{\dd y}\right| = 2\cdot\frac{1}{\sqrt{2\pi}\sigma}e^{-\frac{y}{2}}\cdot\frac{\sigma}{2\sqrt{y}}\\
		&\Rightarrow p(\chi_1^2) = \frac{1}{\sqrt{2\pi}} \cdot \frac{1}{\sqrt{\chi^2}}e^{-\frac{\chi^2}{2}}
	\end{align*} 
	
We note that this diverges for $\chi^2=0$, and that the sum over the solutions gave $2$ for the positive and negative square root. 
	
We now extend the result to $\nu$ degrees of freedom, using the characteristic function $\phi$:
	
	\begin{align*}
		\phi_{x^2}(t) &= \E_x\left[e^{itx^2}\right] = \frac{1}{\sqrt{2\pi}}\int_{-\infty}^{+\infty}e^{itx^2}\cdot e^{-\frac{x^2}{2}}\dd x\\
		&= \frac{1}{\sqrt{2\pi}}\int_{-\infty}^{+\infty}e^{-\left(1-2it\right)\frac{x^2}{2}}\dd x \\
		&\text{Substituting } \sigma= \left(1-2it\right)^{-\frac{1}{2}}\\
		\Rightarrow \phi_{x^2} &=\sigma\underbrace{\int_{-\infty}^{+\infty}\frac{1}{\sqrt{2\pi}\sigma}e^{-\frac{x^2}{2\sigma^2}}\dd x}_{1} = \sigma
	\end{align*}
	
Therefore, we get:
	
	\[
	\phi_{\chi_1^2}(t) = \left(1-2it\right)^{-\frac{1}{2}}
	\]
	
Knowing that the distribution of the sum  is the product of the characteristic functions, we can pass to $\nu$ degrees of freedom by elevating to the power of $\nu$.
	
	\[
	\phi_{\chi_\nu^2}(t) = \left(1-2it\right)^{-\frac{\nu}{2}}
	\]
	
We use the inverse transformation to calculate the pdf of $\chi_\nu^2$:
	
	\[
	p_\nu(\chi^2) = \frac{1}{2\pi}\int_{-\infty}^{+\infty}\phi(t)e^{-itx}\dd t
	\]
	
Using the Gamma function (see \ref{sec:gamma}), through a change of variable, we get equation (\ref{eq:chi_squared}):
	
	\[
	p_\nu(\lambda) = \frac{1}{2^{\frac{\nu}{2}} \Gamma(\frac{\nu}{2})}\lambda^{\frac{\nu}{2}-1}e^{-\frac{\lambda}{2}}
	\]
	
From the derivatives of the characteristic functions we get the moments of this distribution, most notably:
	
	\[
	\E\left[\chi_\nu^2\right] =\nu\quad \Var\left[\chi_\nu^2\right] = 2\nu
	\]
	
Now we know that this distribution was obtained through the sum of $\nu$ variables with finite variance, and from the CLT we can deduce that
	
	\[
	\lim_{\nu \to \infty} \chi_\nu^2 = \N(\nu,2\nu)
	\]


\begin{example}\label{ex:llr_poissonian}
	We have already shown how to compute confidence intervals numerically in the Poisson case. We now derive an approximate solution using the log-likelihood ratio (LLR) method.
	
	Consider a single observation
	\[
	n \sim \text{Poisson}(\mu)
	\]
	The likelihood is
	\[
	p(n;\mu) = \frac{\mu^n e^{-\mu}}{n!}
	\]
	and, from \ref{eq:poisson_mle}, the maximum likelihood estimator is
	\[
	\hat\mu = n
	\]

	We now compute the log-likelihood ratio (LLR):
	\begin{align*}
		\lambda(n, \mu)
		&= 2 \ln \left(\frac{p(n;\hat\mu)}{p(n;\mu)}\right) \\
		&= 2 \left[ \ln p(n; \hat\mu) - \ln p(n;\mu) \right] \\
		&= 2 \left[ \ln \left(\frac{n^n e^{-n}}{n!}\right) - \ln \left(\frac{\mu^n e^{-\mu}}{n!}\right) \right] \\
		&= 2 \left[ n \ln n - n - n \ln \mu + \mu \right] \\
		&= 2 \left[ n \ln \left(\frac{n}{\mu}\right) + \mu - n \right]
	\end{align*}
	
	According to Wilks' theorem, in the asymptotic regime the statistic $\lambda(n,\mu)$ follows a $\chi^2$ distribution with one degree of freedom. Therefore, a confidence interval at confidence level $\mathrm{CL}$ is obtained by solving
	\[
	\lambda(n,\mu) \leq \chi^2_{1,\mathrm{CL}}
	\]
	
	For example, at $\mathrm{CL}=95\%$, the critical value is
	\[
	\chi^2_{1,0.95} \simeq 3.84
	\]
	
	The confidence interval for $\mu$ is thus defined as the set of values satisfying
	\[
	2 \left[ n \ln \left(\frac{n}{\mu}\right) + \mu - n \right] \leq 3.84
	\]
	which must be solved numerically for $\mu$ given the observed $n$.
	
	This procedure produces approximate confidence intervals, often referred to as likelihood-ratio intervals. For each observed value of $n$, one obtains a lower and upper bound $(\mu_\text{low}, \mu_\text{up})$ by finding the roots of
	\[
	\lambda(n,\mu) = \chi^2_{1,\mathrm{CL}}
	\]
	
	The following table compares the likelihood-ratio (Wilks) intervals with the exact Feldman-Cousins construction and the central (equal-tail) intervals, for $n = 0, \dots, 9, 50$:
	
	\input{tables/llr_poisson_intervals.tex}
	
	Differences are most visible at small $n$: the Wilks intervals are typically narrower, while the central ones are more conservative, whereas Feldman-Cousins provides the reference construction with correct coverage. As $n$ increases, all methods converge, reflecting the asymptotic validity of the likelihood-ratio approximation.

	This behavior can be quantified through the coverage function:
	
	\[
	\mathcal{C}(\mu) = \sum_{n \in \text{accepted}} \text{Poisson}(n;\mu)
	\]
	
	where the sum runs over all values of $n$ for which $\mu$ lies inside the corresponding interval. 
	
	As shown in Fig.~\ref{fig:llr_poisson_coverage}, the coverage is not exactly constant and oscillates around the nominal value. These oscillations are a direct consequence of discreteness, as seen before in \ref{fig:binomial_coverage}, and the use of an asymptotic approximation.
	
	\begin{figure}[H]
		\centering
		\input{images/llr_poisson_coverage.pgf}
		\caption{Coverage function $\mathcal{C}(\mu)$ for the LLR-based intervals in the Poisson case. The dashed line represents the nominal confidence level.}
		\label{fig:llr_poisson_coverage}
	\end{figure}
	
	If a more accurate construction is required, one can resort to a fully numerical approach, for instance by computing the distribution of $\lambda(n,\mu)$ via Monte Carlo and inverting it to obtain intervals with exact coverage.
\end{example}

\subsection[A unifying perspective]{A unifying perspective\protect\sidenotemark}
\sidenotetext{The reader is warned that in the following we are going to talk about the core concepts in a mathematically loose way and in a more qualitative manner then previously done. The many imprecision (and possibly blunders) that we are going to make are to be intended as a genuine attempt at making the bigger picture clearer. If you think you already have a good grasp on the matter, you can skip what comes!}
If you're as confused as I was when I first approached the topic of interval estimations (frankly I'm still more than a bit confused :|), it could be helpful to make a real quick and general summary of the concepts introduced. We can exploit the example on the use of LLR on Poissonian (\ref{ex:llr_poissonian}) to sum up some concepts. To build the table we used three different approaches:
\begin{itemize}
	\item Wilks' method
	\item Central interval (equal-tail)
	\item Feldman-Cousins method
\end{itemize}
We can gain some further insight on the matter by better understanding the difference between this methods.

From the point of view of Neyman’s construction, all these approaches aim at defining the same fundamental object: an acceptance region $\mathcal{P}(\theta)$ in the space of possible observations, such that
\[
P(x \in \mathcal{P}(\theta); \theta) = CL
\]

Once this region is specified for every value of $\theta$, the confidence interval corresponding to an observed value $x_ {obs}$ is obtained by inversion:
\[
\text{CI}(x_{\text{obs}}) = \{ \theta: x_{\text{obs}} \in \mathcal{P}(\theta) \}
\]

In this sense, all interval estimation methods differ only in how the acceptance regions $\mathcal{P}(\theta)$ are constructed.

A useful way to classify the methods discussed so far is the following.

\paragraph{Explicit Neyman construction}

The Feldman--Cousins method represents the most direct implementation of Neyman’s idea. For each fixed value of $\theta$, one considers all possible outcomes $x$ and introduces an ordering principle (in this case based on the likelihood ratio)
\[
R(x,\theta) = \frac{p(x;\theta)}{p(x;\hat{\theta}_{MLE})}
\]
The values of $x$ are then included in the acceptance region $\mathcal{P}(\theta)$ in order of decreasing $R(x,\theta)$ until the desired probability content $CL$ is reached:
\[
\int_{x \in \mathcal{P}(\theta)} p(x;\theta)dx = CL
\]

Repeating this procedure for all $\theta$ produces a full mapping
\[
\theta \longrightarrow \mathcal{P}(\theta)
\]
which can be visualized as a confidence belt in the $(\theta, x)$ plane. The interval for a given observation is then obtained by intersecting this belt with the line $x = x_{obs}$.

This construction is conceptually complete and guarantees correct coverage by design, but it is often computationally demanding.

In practice, the acceptance regions are evaluated \emph{numerically} for the different $\theta$, \emph{then} the corresponding interval is extracted from the observed value $x_{obs}$. 

\paragraph{Implicit constructions}

Other methods, such as central intervals and likelihood-ratio (Wilks) intervals, can be understood as implicit implementations of the same idea. In these cases, the acceptance regions are not constructed explicitly for all $\theta$, but are instead obtained by inverting conditions evaluated at the observed value.

For instance:
\begin{itemize}
	\item In the central (equal-tail) construction, the interval is obtained by selecting $\theta$ values such that $x_{obs}$ lies between appropriate quantiles of the distribution.
	\item In the Wilks approach, one considers the log-likelihood ratio
	\[
	\lambda(x,\theta) = 2\ln\left(\frac{p(x;\hat{\theta}_{MLE})}{p(x;\theta)}\right)
	\]
	and defines the interval through the condition
	\[
	\lambda(x_{obs},\theta) \leq q_{CL}
	\]
\end{itemize}

These methods rely on statistics whose distribution is known exactly or asymptotically. In particular, Wilks' theorem provides an asymptotic pivotal quantity, allowing one to avoid the explicit construction of the confidence belt.

In practice, one \emph{directly} evaluates the condition at $x_{obs}$ and solves for $\theta$.

\paragraph{Summary}

All the interval estimation procedures discussed can be viewed as implementations of the same underlying idea:
\begin{itemize}
	\item define a region in the observation space with probability content $CL$
	\item invert this construction to obtain a set of plausible parameter values
\end{itemize}

The difference lies in how this is achieved:
\begin{itemize}
	\item explicitly, by constructing the full Neyman belt (Feldman-Cousins);
	\item implicitly, by exploiting pivotal quantities or approximations (central intervals, Wilks)
\end{itemize}

In all cases, the objective is to obtain intervals with the desired coverage. The choice of method is therefore guided by practical considerations such as computational cost, optimality criteria, and the validity of asymptotic approximations.
\todo{This whole part is really an experiment and I'm not satisfied with the result. I'll try to refine it little by little, because I think that a qualitative overview of the subject can be hugely helpful.}
The difference lies in how this is achieved:
\begin{itemize}
	\item explicitly, by constructing the full Neyman belt (Feldman-Cousins);
	\item implicitly, by exploiting pivotal quantities or approximations (central intervals, Wilks)
\end{itemize}

In all cases, the objective is to obtain intervals with the desired coverage. The choice of method is therefore guided by practical considerations such as computational cost, optimality criteria, and the validity of asymptotic approximations.
\todo{This whole part is really an experiment and I'm not satisfied with the result. I'll try to refine it little by little, because I think that a qualitative overview of the subject can be hugely helpful.}